\documentclass{article}
\usepackage[utf8]{inputenc}
\usepackage{geometry}
\usepackage{graphicx}
\usepackage{pgfplots}
\usepackage{amsmath}
\usepackage{amssymb}
\usepackage{amsfonts}
\usepackage[thmmarks, amsmath, thref]{ntheorem}
\usepackage{mdframed}
\usepackage{amsthm}
\usepackage{physics}
\usepackage{derivative}
\usepackage{hyperref}
\usepackage{wrapfig}
\usepackage{amsthm}
\renewcommand{\theenumi}{\roman{enumi}}

\theoremstyle{nonumberplain}
\theoremheaderfont{\itshape}
\theorembodyfont{\upshape}
\theoremseparator{.}
\theoremsymbol{\ensuremath{\square}}
\theoremsymbol{\ensuremath{\blacksquare}}
\newtheorem{solution}{Solution}
\theoremseparator{. ---}
\theoremsymbol{\mbox{\texttt{;o)}}}

\geometry{a4paper, left=2cm, right=2cm, top=1cm, bottom=2cm}

\pgfplotsset{compat=1.18}
\begin{document}

La ecuación (7) establece que la razón entre los senos de los ángulos de transmisión e incidencia es constante y está dada por la razón de sus velocidades de propagación en cada medio:

\[
\frac{\sin{\theta_T}}{\sin{\theta_I}} = \frac{v_T}{v_I} = \text{constante}
\]

Si se busca determinar la velocidad de la luz en uno de los medios, suponiendo que la del otro medio es conocida con exactitud, podemos reescribir la ecuación como:

\[
v_T = v_I \frac{\sin{\theta_T}}{\sin{\theta_I}}, \quad \quad
v_I = v_T \frac{\sin{\theta_I}}{\sin{\theta_T}}
\]

Dependiendo de qué velocidad es conocida, se despeja la ecuación correspondiente; si se linealiza la ecuación (7) en una representación gráfica de \(\sin{\theta_T}\) vs. \(\sin{\theta_I}\), la pendiente \(m\) de la recta ajustada es:

\[
m = \frac{v_T}{v_I}
\]

Dado que en un experimento típico se obtiene \(m\) a partir del ajuste lineal de datos experimentales, podemos despejar la velocidad desconocida en términos de la otra:

\[
v_T = m v_I, \quad \quad
v_I = \frac{v_T}{m}
\]

según el caso de estudio. Dado que la pendiente \(m\) se obtiene con cierta incertidumbre \(\Delta m\), la propagación de errores se determina mediante:

\[
\Delta v_T = v_I \Delta m,  \quad \quad \Delta v_I = \frac{v_T \Delta m}{m^2}
\]

donde \(\Delta v_T\) y \(\Delta v_I\) representan las incertidumbres en las velocidades de transmisión e incidencia respectivamente. \\


\end{document}
